\section{CRUD Aplikasi Web dengan PHP dan MySQL}

Section ini menyajikan \textbf{studi kasus} CRUD lengkap: Aplikasi Manajemen Produk yang menggabungkan form HTML, PHP, dan MySQLi. Aplikasi terdiri dari beberapa operasi: menampilkan daftar produk (Read), menambah produk baru (Create), mengedit produk (Update), dan menghapus produk (Delete). Seluruh operasi menggunakan prepared statement untuk keamanan \cite{phpManual}.

\begin{contoh}
\textbf{Studi Kasus: Aplikasi CRUD Produk}

Alur aplikasi: (1) Halaman utama menampilkan daftar produk dari database (SELECT); (2) Form ``Tambah Produk'' mengirim data via POST; (3) PHP memvalidasi input, lalu INSERT menggunakan prepared statement; (4) Tombol ``Edit'' memuat data produk ke form, lalu UPDATE; (5) Tombol ``Hapus'' mengirim ID via POST, lalu DELETE. Semua output menggunakan \code{htmlspecialchars()} untuk mencegah XSS.
\end{contoh}

\begin{webcode}[language=PHP, caption={CRUD Produk - Tampil dan Tambah (Read + Create)}]
<?php
require_once 'config.php'; // koneksi $conn

// CREATE - Insert produk baru
if ($_SERVER['REQUEST_METHOD'] === 'POST' && ($_POST['aksi'] ?? '') === 'tambah') {
  $nama  = trim($_POST['nama'] ?? '');
  $harga = floatval($_POST['harga'] ?? 0);

  if (!empty($nama) && $harga > 0) {
    $stmt = $conn->prepare("INSERT INTO produk (nama, harga) VALUES (?, ?)");
    $stmt->bind_param("sd", $nama, $harga);
    $stmt->execute();
    $stmt->close();
    header("Location: produk.php");
    exit;
  }
}

// READ - Ambil semua produk
$result = $conn->query("SELECT id, nama, harga FROM produk ORDER BY id DESC");
$produkList = $result->fetch_all(MYSQLI_ASSOC);
$result->free();
?>
<!DOCTYPE html>
<html lang="id">
<head><title>CRUD Produk</title></head>
<body>
<h2>Daftar Produk</h2>
<table border="1" cellpadding="5">
  <tr><th>ID</th><th>Nama</th><th>Harga</th><th>Aksi</th></tr>
  <?php foreach ($produkList as $p): ?>
  <tr>
    <td><?= $p['id'] ?></td>
    <td><?= htmlspecialchars($p['nama']) ?></td>
    <td>Rp <?= number_format($p['harga'], 0, ',', '.') ?></td>
    <td>
      <a href="edit.php?id=<?= $p['id'] ?>">Edit</a>
      <form method="POST" style="display:inline">
        <input type="hidden" name="aksi" value="hapus">
        <input type="hidden" name="id" value="<?= $p['id'] ?>">
        <button onclick="return confirm('Hapus?')">Hapus</button>
      </form>
    </td>
  </tr>
  <?php endforeach; ?>
</table>

<h3>Tambah Produk</h3>
<form method="POST">
  <input type="hidden" name="aksi" value="tambah">
  <label>Nama: <input type="text" name="nama" required></label><br>
  <label>Harga: <input type="number" name="harga" min="0" required></label><br>
  <button type="submit">Tambah</button>
</form>
</body>
</html>
\end{webcode}

\begin{webcode}[language=PHP, caption={CRUD Produk - Update dan Delete}]
<?php
require_once 'config.php';

// DELETE
if ($_SERVER['REQUEST_METHOD'] === 'POST' && ($_POST['aksi'] ?? '') === 'hapus') {
  $id = intval($_POST['id'] ?? 0);
  if ($id > 0) {
    $stmt = $conn->prepare("DELETE FROM produk WHERE id = ?");
    $stmt->bind_param("i", $id);
    $stmt->execute();
    $stmt->close();
  }
  header("Location: produk.php");
  exit;
}

// UPDATE
if ($_SERVER['REQUEST_METHOD'] === 'POST' && ($_POST['aksi'] ?? '') === 'update') {
  $id    = intval($_POST['id'] ?? 0);
  $nama  = trim($_POST['nama'] ?? '');
  $harga = floatval($_POST['harga'] ?? 0);

  if ($id > 0 && !empty($nama) && $harga > 0) {
    $stmt = $conn->prepare("UPDATE produk SET nama = ?, harga = ? WHERE id = ?");
    $stmt->bind_param("sdi", $nama, $harga, $id);
    $stmt->execute();
    $stmt->close();
  }
  header("Location: produk.php");
  exit;
}
?>
\end{webcode}

